\documentclass[border=15pt]{standalone} % 去掉 tikz 选项，完美贴合边缘
\usepackage{ctex}      % 提供中文支持
\usepackage{array}
\usepackage{tabularx}
\usepackage{booktabs}  % 优雅的三线表
\usepackage{xcolor}
\usepackage{colortbl}

% 保持与前面表格完全一致的深灰蓝表头色
\definecolor{headerbg}{RGB}{52, 73, 94} 
\definecolor{headertext}{RGB}{255, 255, 255} 

\begin{document}

% 设置无衬线字体
\sffamily
% 行高设为 1.5，保持舒适的阅读留白
\renewcommand{\arraystretch}{1.5}

% 总宽度设为 16cm，第一列加粗占 4.5cm（为了放下 RHEL 那么长的名字），第二列自动铺满
\begin{tabularx}{16cm}{>{\bfseries}p{4.5cm} X}
	\rowcolor{headerbg}
	\toprule
	\textcolor{headertext}{\textbf{核心发行版}} & \textcolor{headertext}{\textbf{核心特点与生态地位}}                                 \\
	\midrule

	Debian                                 & 1993 年发布，以系统稳定性与包管理的严谨性著称。广泛应用于服务器领域，为 Ubuntu 等众多衍生发行版的底层基石。               \\
	\addlinespace[1ex]

	Red Hat Enterprise Linux (RHEL)        & 1995 年发布的商业级发行版。面向企业级服务器市场，提供长期技术支持验证与高可用性架构。                              \\
	\addlinespace[1ex]

	Fedora                                 & Red Hat 赞助的社区版本。定位为新技术的上游测试平台，更新策略激进，适合开发者与验证最新 Linux 特性。                  \\
	\addlinespace[1ex]

	Slackware                              & 1993 年发布，严格遵循传统 Unix 设计哲学。系统配置高度依赖手动编辑文本文件，无自动依赖解析的包管理器。                   \\
	\addlinespace[1ex]

	openSUSE                               & 由 SUSE 公司赞助的社区发行版。提供稳定的系统环境，并内置著名的 YaST (Yet another Setup Tool) 综合配置管理工具。 \\

	\bottomrule
\end{tabularx}

\end{document}